\section{Evaluation Method}

\subsection{Scope and study design}

We conduct a software-system evaluation of the interaction contract, its
runtime enforcement, and its authoring controls. The evaluation does not score
generated answers or human experience. It has three parts: (1) an exhaustive
structural routing benchmark over the authored case corpus, (2) contract tests
that execute valid and deliberately broken Unity/backend traces, and (3)
transaction tests for the placement-authoring loop.

The structural benchmark compares two executable adapters. The
\emph{direct-wiring} adapter reads the scene-semantics, agent-role, and handoff
artifacts used by the existing implementation. The
\emph{model-grounded} adapter reads a fresh native \systemname{} instance
regenerated in memory from the corresponding version-0.5 case model. Both
adapters expose the same normalized operation: resolve the owner of an object
and classify a route from a given starting agent as local, directly permitted,
transitively reachable, or unreachable. This is an adapter-level structural
comparison, not a comparison of conversational models or Unity frame rates.

\subsection{Cases and fixed generation inputs}

The corpus contains a career-fair room, a classroom with a dinosaur exhibit,
and a food or trade-fair room. They provide 27, 36, and 30 routable scene
objects and 5, 4, and 6 agents, respectively. The object universe is the union
of the agents' explicitly grounded object identifiers. One zone-only semantic
reference in the food/trade-fair case is disclosed and excluded because it is
not an executable agent-grounding target.

Scene semantics, agent roles, handoff matrices, and knowledge entries had been
proposed in earlier pipeline runs using
\texttt{gpt-4.1-mini-2025-04-14} at temperature 0.2. We treat those stored
outputs as fixed case inputs. The final model assembly, comparison, mutations,
metrics, and tables make no model or network calls.

\subsection{Responsibility and routing probes}

Both adapters apply the same precedence rule: an explicit asset assignment
outranks a matching object-group assignment, which outranks a matching zone.
More than one owner at the highest active tier is ambiguous and must fail
closed. For every routable object we first compare the resolved owner between
adapters. We then pair that object with every possible starting agent in its
case and compare the route classification. The resulting corpus is exhaustive
for these three fixed models; it is not sampled from a larger population.

Natural inputs assign every routable object at the asset tier. We therefore
keep fallback tests separate from natural-corpus results. For one copied object
per case, we derive three synthetic probes: remove the asset assignment to
exercise group fallback, also remove the group assignment to exercise zone
fallback, and add a competing owner at the highest active tier to exercise
ambiguity rejection. The same normalized candidate sets are supplied to both
adapters. Because the direct-wiring source has no native object-group
relationship, these probes test the decision rule, not equivalent native
authoring expressiveness.

\subsection{Fault injection and edit counts}

The common mutation denominator contains three semantically equivalent changes
per case: remove ownership of an object, assign a second owner, and introduce a
dangling handoff. For each adapter and mutation we start from an accepted
copy, apply a deterministic patch, and record whether a new validation error is
raised and whether it identifies the affected relationship. We also count
edited top-level artifacts and changed references. These counts describe the
scripted representation patches; they are not elapsed developer effort.

A separate V2-only suite removes a capability provider, duplicates a source
agent identifier, or reassigns a domain capability to the runtime
orchestrator. We report the canonical model validator and the stricter
pipeline provider contract separately, because combining them would obscure
which rule detected a fault. Expected-valid copies measure false positives.

\subsection{Local structural runtime}

For descriptive overhead, each adapter repeatedly constructs its normalized
responsibility view and enumerates all object-by-start-agent routes. After four
warm-ups, we alternate adapter order for 40 repetitions per case and measure
with a monotonic nanosecond clock. JSON input/output, version-0.5-to-V2
generation, Unity, and language-model latency are excluded. We report medians
and the V2/direct ratio, but do not infer a population-level performance
advantage.

\subsection{Runtime contract tests}

Backend tests pin a V2 model at session setup and exercise deictic and
non-deictic chat requests. Negative cases include a missing spatial context,
an unknown entity, an obsolete model hash, an ambiguous target, an invalid
active agent, and an unmodeled handoff. Each test snapshots session history and
checks that a rejected request caused no mutation or model call. Positive cases
check the server-derived entity, group, zone, responsible agent, evidence
references, routing reason, and modeled binding returned to Unity.

Unity batch-mode smokes execute the same interaction-evidence component used by
the client. A valid target and route must reach completion; a mismatched
response entity, missing deictic context, ambiguous selection, or forbidden
handoff must produce a failed assertion and block answer completion.
Non-deictic interactions are represented explicitly and cannot claim spatial
grounding evidence.

\subsection{Placement-authoring transaction}

The editor's authoring test begins from a validated FunctionalMLDS draft. It
inspects the current revision, previews a structured placement change without
writing files, applies the change through the deterministic placement and model
stages, and then either accepts or discards it. A subsequent undo restores the
last accepted snapshot without restarting the backend. Tests compare
before/after values and artifact bytes, reject stale revisions and unknown
agents, force a regeneration failure to exercise rollback, and ensure that a
runtime project cannot be committed while a change awaits acceptance.
Responsibility editing is not evaluated or presented as implemented.

\subsection{Measures, analysis, and reproducibility}

For RQ1 we report unique, ambiguous, and unassigned ownership; route classes;
one-hop coverage; graph reachability; adapter parity; and the separate
fallback probes. For RQ2 we report mutation detection, false positives,
localization, representation edits, and descriptive local runtime. For RQ3 we
report pass/fail outcomes for request validation, scenario evidence, rollback,
discard, and undo.

All proportions include their numerator and denominator. We do not attach
confidence intervals or significance tests to exhaustive deterministic probes
over three purposively authored cases: such intervals would imply a sampling
model the design does not support. Repetitions characterize local timing noise
only. A single API-free command records software versions, source and input
hashes, immutable JSON results, and generated tables. The reported evaluation
contains no human participants or personal data and supports no claims about
usability, preference, trust, or workload.
